\begin{table}[htbp]
\centering
\footnotesize
\caption{\textbf{1.3} Bit\'acora de limpieza. El orden en que se aplican las reglas es parte de la bit\'acora.}
\label{tab:p1_bitacora}
\begin{tabular}{clrr}
\toprule
\textbf{orden} & \textbf{regla} & \textbf{filas eliminadas} & \textbf{monto (\pounds)} \\
\midrule
1 & R1 duplicados exactos & 5\,268 & 21\,740.98 \\
2 & R2 facturas con prefijo distinto de digito o C & 3 & -11\,062.06 \\
3 & R3 StockCode que no es venta al cliente & 149 & -239\,669.22 \\
4 & R4 precio unitario <= 0 & 2\,510 & 0.00 \\
5 & R5 cantidad negativa sin factura de cancelacion & 0 & 0.00 \\
6 & R6 filas sin CustomerID (SE QUEDAN) & 0 & 0.00 \\
7 & TOTAL & 7\,930 & -228\,990.30 \\
\bottomrule
\end{tabular}
\\[0.3em]{\scriptsize \textbf{Qu\'e leer:} R5 elimina cero filas porque R4 ya se llev\'o las 1\,336 filas de cantidad negativa, que ten\'ian todas precio cero. La tabla siguiente muestra qu\'e pasa si se invierte el orden.}
\end{table}
